% !TeX program = lualatex
% Lebenslauf – Erlind Sejdiu
% Kompilieren mit:  lualatex lebenslauf.tex

\documentclass[a4paper,10pt]{article}

% ─── Pakete ──────────────────────────────────────────────────────────────────
\usepackage{fontspec}
\usepackage[ngerman]{babel}
\usepackage[left=2.1cm, right=2.1cm, top=1.4cm, bottom=1.0cm]{geometry}
\usepackage{microtype}
\usepackage{enumitem}
\usepackage{tabularx}
\usepackage{xcolor}
\usepackage{titlesec}
\usepackage[hidelinks]{hyperref}

% ─── Schrift ─────────────────────────────────────────────────────────────────
% TeX Gyre Pagella per Dateinamen laden (funktioniert plattformübergreifend in TeX Live & MiKTeX)
\setmainfont{texgyrepagella}[
    Extension = .otf,
    UprightFont = *-regular,
    BoldFont = *-bold,
    ItalicFont = *-italic,
    BoldItalicFont = *-bolditalic
]

% ─── Farben: nur Schwarz und zwei Grautöne ───────────────────────────────────
\definecolor{ink}{HTML}{1C1C1C}
\definecolor{muted}{HTML}{6E6E6E}
\definecolor{hairline}{HTML}{C4C4C4}
\color{ink}

\hypersetup{pdfauthor={Erlind Sejdiu}, pdftitle={Lebenslauf -- Erlind Sejdiu}}

% ─── Überschriften: Kapitälchen statt Versalien und Farbbalken ───────────────
\titleformat{\section}
    {\LARGE\scshape\color{ink}}{}{0em}{}
    [\vspace{-0.65em}\textcolor{hairline}{\rule{\linewidth}{0.4pt}}]
\titlespacing*{\section}{0pt}{0.7em}{0.4em}

\setlength{\parindent}{0pt}
\setlength{\parskip}{0pt}

% ─── Breite der Datumsspalte ─────────────────────────────────────────────────
\newlength{\datecol}\setlength{\datecol}{3.8cm}

% \eintrag{Zeitraum}{Position}{Institution}
\newcommand{\eintrag}[3]{%
    \noindent
    \begin{tabularx}{\textwidth}{@{}p{\datecol}@{}X@{}}
        \textcolor{muted}{#1} & \textbf{#2}\ifx\relax#3\relax\else, #3\fi
    \end{tabularx}%
    \par\nobreak
}

% \block{Zeitraum}{Titel}{Fließtext} – falls kein Zeitraum angegeben, volle Breite
\newcommand{\block}[3]{%
    \noindent
    \ifx\relax#1\relax
        \textbf{#2}\par #3\par\vspace{0.35em}%
    \else
        \begin{tabularx}{\textwidth}{@{}p{\datecol}@{}X@{}}
            \textcolor{muted}{#1} & \textbf{#2}\par #3
        \end{tabularx}\vspace{0.2em}%
    \fi
}

% Aufzählung unterhalb eines Eintrags, bündig zur rechten Spalte
\newenvironment{details}
    {\begin{itemize}[leftmargin=\dimexpr\datecol+0.9em\relax, labelsep=0.5em,
                     labelindent=\datecol, topsep=0.15em, itemsep=0.1em,
                     parsep=0pt, label=\textcolor{hairline}{\textbf{\textperiodcentered}}]}
    {\end{itemize}\vspace{0.2em}}

\pagestyle{empty}

% =============================================================================
\begin{document}

% ─── Kopf ────────────────────────────────────────────────────────────────────
{\huge Erlind Sejdiu}\\[0.35em]
{\small\color{muted}
Blücherstraße 18, Karlsruhe \enspace\textbar\enspace
0157\,3547\,4746 \enspace\textbar\enspace
\href{mailto:erlindamprogrammieren@gmail.com}{erlindamprogrammieren@gmail.com}}\\[0.45em]
\textcolor{hairline}{\rule{\linewidth}{0.7pt}}

\vspace{0.5em}

% ─── Bildung ─────────────────────────────────────────────────────────────────
\section{Bildung}

\eintrag{seit 3/2026}{M.\,Sc.\ Informatik}{Hochschule Karlsruhe}
\begin{details}
    \item Schwerpunkt Software Engineering, Notendurchschnitt nach dem 1.~Semester 1,0
\end{details}

\eintrag{10/2022 -- 02/2026}{B.\,Sc.\ Technische Informatik}{Hochschule Pforzheim \normalfont(Abschlussnote 1,3)}
\begin{details}
    \item Bachelorarbeit bei der ProSeS GmbH (Note 1,0): \glqq Refactoring des Benachrichtigungs\-moduls
          zur Optimierung eines Betriebsdatenerfassungssystems\grqq
\end{details}

% ─── Berufserfahrung ─────────────────────────────────────────────────────────
\section{Berufserfahrung}

% TODO: genauen Monat des Beginns eintragen
\eintrag{seit 09/2025}{Werkstudent}{ProSeS GmbH}
\begin{details}
    \item Weiterentwicklung eines Systems zur Betriebsdatenerfassung
    \item Refactoring des Benachrichtigungsmoduls, Absicherung durch Unit-Tests
\end{details}

\eintrag{09/2024 -- 02/2025}{Praxissemester, R\&D Embedded Systems}{Vector Informatik GmbH}
\begin{details}
    \item Evaluation von Technologien und Standards im Automotive-Umfeld
    \item Aufbau und Programmierung eines Proof of Concept auf einem Automotive-Switch
\end{details}

\eintrag{07/2023 -- 02/2024}{Werkstudent Cybersecurity}{Codewerk GmbH}
\begin{details}
    \item Entwicklung eines GUI-Fuzzers auf Basis von Radamsa
    \item Auswertung in Kibana-Dashboards, Containerisierung mit Docker Compose
\end{details}

\eintrag{02/2023 -- 06/2023,\newline 02/2025 -- 08/2025}{Wissenschaftliche Hilfskraft}{Hochschule Pforzheim\newline \normalfont Tutor für Softwareentwicklung und Software Engineering}
\begin{details}
    \item Betreuung der Labore Softwareentwicklung und Software Engineering\,2
    \item Unterstützung bei Programmieraufgaben und beim Schreiben von Software-Tests
\end{details}

% ─── Projekte ────────────────────────────────────────────────────────────────
\section{Projekte}

\block{}{Aktienmarkt-Analysetool}
    {Simulation von Kursverläufen und automatisiertes Auffinden geeigneter Handelsalgorithmen.}

\block{}{Projekt-Management Applikation}
    {Dediziertes Werkzeug zur Planung des Alltags und paralleler Projekte, exakt auf persönliche Anforderungen zugeschnitten.}

\block{}{Movelink -- Bewegungsanalyse-Tool}
    {Ein IMU-basiertes Analysetool zur Erkennung und Vorhersage von Bewegungsmustern.}

% ─── Engagement ──────────────────────────────────────────────────────────────
\section{Ehrenamtliches Engagement}

\block{2023 -- 2025}{Fachschaft Technik, Hochschule Pforzheim}
    {Organisation und Planung von Veranstaltungen für Studierende der Fakultät Technik.}

\block{2023 -- 2024}{Studierendenrat, Hochschule Pforzheim}
    {Vertretung studentischer Interessen, Mitorganisation hochschulweiter Veranstaltungen.}

% ─── Kenntnisse ──────────────────────────────────────────────────────────────
\section{Kenntnisse}

\begin{tabularx}{\textwidth}{@{}p{\datecol}@{}X@{}}
    \textcolor{muted}{Programmierung} & C++, C, Python, C\#, MATLAB, Bash, VHDL \\[0.15em]
    \textcolor{muted}{Werkzeuge}      & Git, CMake, Make, GoogleTest, Docker, Visual Studio, LTspice \\[0.15em]
    \textcolor{muted}{Sprachen}       & Deutsch (Muttersprache), Englisch (verhandlungssicher) \\
\end{tabularx}

% ─── Ort und Datum ───────────────────────────────────────────────────────────
\vfill

\end{document}
